\section{Giới thiệu}

\subsection{Bối cảnh đề tài}
Trong những năm gần đây, xu hướng nuôi thú cưng tại Việt Nam, đặc biệt ở các đô thị lớn như Thành phố Hồ Chí Minh, đã có bước chuyển dịch mạnh mẽ từ việc nuôi để giữ nhà sang xem thú cưng như những thành viên thân thiết trong gia đình (hiện tượng nhân hóa thú cưng -- pet humanization). Theo các khảo sát thị trường đô thị, nhu cầu sử dụng các dịch vụ chăm sóc thú cưng định kỳ như tắm gội chuyên sâu, vệ sinh tai móng, cắt tỉa tạo kiểu lông (grooming) và dịch vụ trông giữ (pet hotel) ngày càng gia tăng vượt bậc. 

Tuy nhiên, phần lớn các cơ sở dịch vụ thú cưng hiện nay vẫn vận hành dựa trên các phương thức truyền thống mang tính thủ công cao: tiếp nhận đặt lịch qua các cuộc gọi thoại trực tiếp, tin nhắn mạng xã hội (Facebook Fanpage, Zalo cá nhân) hoặc ghi chép tay trên sổ quản lý tiếp nhận tại quầy lễ tân. Trong khi đó, nhóm đối tượng khách hàng chủ đạo sử dụng dịch vụ này lại là những người trẻ, nhân viên văn phòng bận rộn với lịch trình làm việc hành chính dày đặc, ít có thời gian chờ đợi hoặc theo dõi sát sao tại cửa hàng. Sự mất cân đối giữa phương thức tương tác thủ công truyền thống của các cơ sở chăm sóc với thói quen sử dụng công nghệ số tiện lợi, tức thì của chủ nuôi hiện đại đã tạo ra nhiều rào cản và trải nghiệm không mong muốn trong quá trình sử dụng dịch vụ.

Đề tài \textbf{``Hệ thống hỗ trợ đặt lịch, gửi yêu cầu và theo dõi quá trình chăm sóc thú cưng''} được nghiên cứu và phát triển nhằm ứng dụng các nguyên lý Tương tác Người -- Máy (HCI) hiện đại~\cite{Norman2013, Nielsen1994}, số hóa toàn diện quy trình giao tiếp hai chiều giữa chủ nuôi và cơ sở cung cấp dịch vụ, từ đó mang lại trải nghiệm tiện lợi, an tâm và minh bạch tuyệt đối cho người dùng cuối.

\subsection{Phát biểu bài toán}
Thông qua tài liệu Proposal đã được phê duyệt và quá trình khảo sát người dùng thực tế, bài toán cốt lõi mà đồ án tập trung giải quyết xuất phát từ các bất cập nghiêm trọng trong quy trình chăm sóc thú cưng thủ công hiện tại:
\begin{itemize}
    \item \textbf{Khó khăn khi đặt lịch hẹn}: Chủ nuôi phải liên hệ qua điện thoại hoặc nhắn tin nhiều lần để hỏi lịch trống, phải chờ đợi nhân viên rà soát lịch hẹn thủ công. Điều này khiến người dùng rơi vào thế bị động, dễ bị trùng lịch hoặc lỡ mất khung giờ mong muốn trong các dịp cao điểm cuối tuần.
    \item \textbf{Xác nhận thiếu tin cậy và dễ sai lệch}: Việc xác nhận lịch hẹn qua tin nhắn văn bản không đồng bộ tự động vào hệ thống quản lý, dẫn đến rủi ro nhân viên quên ghi sổ hoặc ghi nhận nhầm lẫn thông tin thời gian, loại dịch vụ và nhân sự phụ trách.
    \item \textbf{Thất lạc dặn dò y tế và yêu cầu đặc biệt}: Thú cưng thường có tiền sử sức khỏe riêng biệt (như dị ứng với một số thành phần xà phòng thơm, bệnh da liễu nhạy cảm, vừa phẫu thuật cần tránh nước, hoặc tính cách hoảng sợ khi sấy lông). Trong quy trình cũ, chủ nuôi phải liên tục nhắc đi nhắc lại các lưu ý này mỗi khi gửi thú cưng. Nhân viên tiếp nhận tại quầy ghi chép sơ sài hoặc không bàn giao đầy đủ cho kỹ thuật viên trực tiếp tắm cắt, dẫn đến tai nạn tái phát dị ứng nguy hiểm hoặc thú cưng bị hoảng loạn.
    \item \textbf{Thiếu cập nhật tiến độ theo thời gian thực}: Trong suốt khoảng thời gian 2--3 giờ thú cưng được gửi tại cơ sở, chủ nuôi hoàn toàn không nắm bắt được thú cưng của mình đang ở giai đoạn nào (đã được tắm chưa, đang cắt tỉa hay đã xong). Trạng thái ``mù mờ'' thông tin này gây nên tâm lý bất an, lo lắng cho chủ nuôi nhưng họ lại ngại gọi điện làm phiền nhân viên đang bận tay thao tác.
    \item \textbf{Thiếu lịch sử chăm sóc có cấu trúc}: Toàn bộ dữ liệu về các lần sử dụng dịch vụ trước (loại dầu tắm đã dùng có hợp da hay không, độ dài cắt tỉa phù hợp, tình trạng da lông phát hiện lúc tắm, chi phí dịch vụ) đều bị trôi mất trong các chuỗi tin nhắn trò chuyện rời rạc, gây khó khăn cho việc theo dõi sức khỏe lâu dài của thú cưng.
\end{itemize}

Hệ quả là cả chủ nuôi lẫn cơ sở chăm sóc đều lãng phí nhiều thời gian liên lạc lặp đi lặp lại, gia tăng nguy cơ sai sót y tế và làm suy giảm sự gắn kết tin cậy giữa hai bên.

\subsection{Mục tiêu đề tài}
Đồ án hướng tới các mục tiêu cụ thể, bám sát định hướng giải pháp đã cam kết trong Proposal:
\begin{itemize}
    \item \textbf{Mục tiêu tương tác \& Trải nghiệm người dùng (UX/Usability Goals)}:
    \begin{enumerate}
        \item \textit{Tính hiệu quả (Effectiveness)}: Đảm bảo chủ nuôi đặt lịch chính xác 100\% vào các khung giờ trống thực tế và thông tin dặn dò y tế được chuyển giao nguyên vẹn tới kỹ thuật viên phụ trách.
        \item \textit{Hiệu suất (Efficiency)}: Rút ngắn thời gian thao tác đặt lịch xuống dưới 1 phút; cung cấp khả năng tra cứu tiến độ và lịch sử chỉ với 1--2 thao tác chạm trên thiết bị di động.
        \item \textit{Phòng ngừa sai sót (Error Prevention)}: Tự động khóa các khung giờ đã kín; thiết kế cơ chế cảnh báo thị giác nổi bật (visual alert) đối với các trường hợp dị ứng hoặc chống chỉ định y tế.
        \item \textit{Sự an tâm \& Hài lòng (Satisfaction \& Peace of Mind)}: Xóa bỏ hoàn toàn trạng thái lo âu của chủ nuôi khi gửi thú cưng bằng hệ thống cập nhật 4 mốc thời gian thực minh bạch.
    \end{enumerate}
    \item \textbf{Mục tiêu sản phẩm thiết kế}:
    \begin{enumerate}
        \item Xây dựng hoàn chỉnh kiến trúc thông tin, luồng tác vụ và bộ khung giao diện Wireframe chuẩn Mobile-first.
        \item Thiết kế bản mẫu tương tác độ nét cao (High-fidelity Interactive Prototype) trên Figma thể hiện trọn vẹn 4 mục tiêu cốt lõi của Persona chính (Persona 1).
        \item Đánh giá tính khả dụng thực tế thông qua các chỉ số định lượng (Task Completion Rate, Time on Task, Error Rate, Điểm khảo sát Likert/SUS) để làm cơ sở tinh chỉnh thiết kế cuối cùng.
    \end{enumerate}
\end{itemize}

\subsection{Phạm vi đề tài}
\begin{itemize}
    \item \textbf{Phạm vi thực hiện (In-Scope)}:
    \begin{itemize}
        \item Tập trung toàn diện vào trải nghiệm giao diện người dùng của \textbf{Chủ nuôi thú cưng (Pet Owners)} trên nền tảng Web di động (Mobile Web).
        \item Triển khai trọn vẹn \textbf{4 trụ cột chức năng cốt lõi}:
        \begin{enumerate}
            \item \textit{Đặt lịch có xác nhận tức thì}: Chọn thú cưng, chọn gói dịch vụ, duyệt khung giờ trống và nhận thông báo xác nhận tự động.
            \item \textit{Hồ sơ thú cưng \& Dặn dò đặc biệt}: Quản lý thông tin sức khỏe, lưu ý dị ứng/thuốc và tự động đồng bộ vào đơn tiếp nhận dịch vụ kèm mã QR Check-in.
            \item \textit{Theo dõi tiến độ thời gian thực (Live Tracking)}: Thanh tiến trình hiển thị minh bạch 4 mốc trạng thái chuẩn (\textit{Đã nhận thú cưng} $\rightarrow$ \textit{Đang chăm sóc} $\rightarrow$ \textit{Hoàn tất} $\rightarrow$ \textit{Chờ đón}) kết hợp thông báo đẩy (Push Notification).
            \item \textit{Lịch sử chăm sóc cá nhân hóa}: Lưu trữ nhật ký từng buổi dịch vụ, tên sản phẩm chuyên dụng đã sử dụng, ghi chú kỹ thuật viên và hóa đơn điện tử minh bạch.
        \end{enumerate}
    \end{itemize}
    \item \textbf{Phạm vi không thực hiện (Out-of-Scope)}:
    \begin{itemize}
        \item Không xây dựng toàn bộ hệ thống quản trị nội bộ ERP phức tạp cho chuỗi cửa hàng (quản lý kho bãi lớn, quản lý chấm công, tính lương nhân sự).
        \item Không tích hợp luồng phát trực tiếp camera giám sát vật lý tại chuồng do giới hạn bảo mật và băng thông hạ tầng mạng trong phạm vi đồ án môn học.
    \end{itemize}
\end{itemize}
